\chapter{二次量子化}
\renewcommand{\currentchapterpath}{chapters/ch02/}

\begin{flushright}
\textit{“名不正，则言不顺；言不顺，则事不成。” ——《论语·子路》
\footnote{合适的语言是讨论问题的前提。二次量子化正是多体物理的“名”与“言”。}}
\end{flushright}

二次量子化（second quantization）把多体问题表述在 Fock 空间中：粒子数可以变化，交换对称性由算符对易/反对易关系自动实现。本章从量子力学算符复习出发，讨论费米/玻色全同粒子与产生湮灭算符，按 Giuliani--Vignale 给出 jellium、对关联与 Wigner 晶体，并按 Gross 补充密度算符、声子简正模、模型层次、相干态与电子–声子耦合。


\chapterinput{operators}
\chapterinput{identical}
\chapterinput{bosons}
\chapterinput{creation}
\chapterinput{unitary}
\chapterinput{density_ops}
\chapterinput{jellium}
\chapterinput{pair_g}
\chapterinput{thermo}
\chapterinput{wigner}
\chapterinput{models}
\chapterinput{phonons_gross}
\chapterinput{coherent}
\chapterinput{elph}
